\scanpage{11}
\mspage{4}

The group $\mathbb Z/6\mathbb Z$ acts transitively on $({}_2X_S)\setminus\{0\}$ \unsure{(}for $\pm1$ acts trivially - thus \({}_2X_S\setminus\{0\}\) is a torsor under
\[
(\mathbb Z/6\mathbb Z)/\pm1\simeq\mathbb Z/3\mathbb Z
\]
\deleted{(\unsure{conjugate}} (\unsure{corresponding|canonical} generator $\unsure{\simeq}\,-\bar\zeta\,\unsure{j|\zeta}=\exp\frac{2\pi}{6}$),
\deleted{\unsure{in this case}} \inserted{the torsor} is identified \unsure{:|with} $S/\pm1$, by associating to each $s\in S$ the midpoint of the radius $[0,s]$, or rather its image \inserted{)} in $X_S$.

The set of fixed points of $\mathbb Z/3\mathbb Z\inserted{\subset\mathbb Z/6\mathbb Z}$ on $X_S$ has cardinality 3, it is obtained by taking the barycentres of the six triangles of the central subdivision (\unsure{bibarycentric}) of the hexagon (circled points) (such as $(1-\bar\zeta)/3$ and so on in the pinned case), they group into two \unsure{packets|regiments} of 3 \unsure{mutually} congruent \unsure{under|modulo} $\Gamma$, \unsure{giving|give} two elements (symmetric) of $X_S$, of order 3,

\margin{10}
